\documentclass{article}
\usepackage{fontspec}
\usepackage{xunicode}
\usepackage{polyglossia}
\setmainlanguage{french}
\usepackage{enumitem}
\usepackage{amssymb}
\usepackage{hyperref}

\begin{document}

% ---- liste imbriquée ----
Les pieuvres sont des créatures d'une intelligence déconcertante.
\begin{itemize}
\item elles peuvent ouvrir un bocal fermé depuis l'intérieur ;
\item elles utilisent des outils :
	\begin{enumerate}
	\item transporter des coquilles de noix de coco pour s'abriter
	\item empiler des rochers pour barricader leur tanière
	\end{enumerate}
\item elles reconnaissent les visages humains individuellement.
\end{itemize}
Aucun vertébré n'a le monopole de l'ingéniosité.

% ---- liste avec des carrés ----
Avant de partir en randonnée, je prends~:
\begin{itemize}[label=$\square$]
\item carte et boussole ;
\item trousse de premiers secours ;
\item ration de secours (barres énergétiques, fruits secs).
\end{itemize}
La documentation du package enumitem, qui permet de
personnaliser ces marqueurs, est disponible sur
\url{https://ctan.org/pkg/enumitem}. Je veux ici une liste qui prend le carré comme puce.

% ---- liste avec alpha ----
Plan du chapitre sur la première guerre mondiale~:
\begin{enumerate}[label=\Alph*)]
\item L’entrée en guerre et la formation des fronts ;
\item Les opérations sur le front Ouest en 1915 et en 1916 ;
\item La guerre totale ;
\end{enumerate}

% ---- liste en deux parties (avec resume) ----
Recette de la tarte aux pommes~:
\begin{enumerate}[label=\arabic*)]
\item mélanger la farine, le beurre froid et une pincée de sel ;
\item ajouter un peu d'eau froide et former une boule sans trop pétrir ;
\item laisser reposer la pâte au réfrigérateur pendant trente minutes.
\end{enumerate}
Pendant ce temps, on prépare la garniture~:
\begin{enumerate}[label=\arabic*),resume]
\item éplucher les pommes, les couper en fines lamelles ;
\item les disposer en rosace sur le fond de tarte précuit.
\end{enumerate}

\end{document}